\documentclass[11pt]{article}
\usepackage[a4paper,margin=2.2cm]{geometry}
\usepackage[T1]{fontenc}
\usepackage[utf8]{inputenc}
\usepackage[french]{babel}
\usepackage{lmodern}
\usepackage{amssymb}
\usepackage{microtype}
\usepackage{xcolor}
\usepackage{titlesec}
\usepackage{enumitem}
\usepackage{fancyhdr}
\usepackage{listings}
\usepackage{tcolorbox}
\tcbuselibrary{skins,breakable}

\definecolor{bcg}{HTML}{177B4C}
\definecolor{bg}{HTML}{F4F6F5}
\definecolor{kw}{HTML}{0B6E99}

\titleformat{\section}{\Large\bfseries\color{bcg}}{\thesection}{0.6em}{}
\titleformat{\subsection}{\large\bfseries}{\thesubsection}{0.6em}{}
\setlist{leftmargin=1.4em}
\pagestyle{fancy}\fancyhf{}
\rhead{\small BCG X CodeSignal Prep}\lhead{\small Support client}
\cfoot{\thepage}

\lstset{
  basicstyle=\ttfamily\small, backgroundcolor=\color{bg},
  frame=single, framesep=5pt, rulecolor=\color{bg!60!black},
  breaklines=true, keywordstyle=\color{kw}\bfseries,
  language=Python, showstringspaces=false, columns=fullflexible,
}

\newtcolorbox{pat}{colback=bcg!8, colframe=bcg, title=SCH\'EMA \`A RECONNA\^ITRE,
  fonttitle=\bfseries, breakable, boxrule=0.6pt}
\newtcolorbox{trap}{colback=red!5, colframe=red!60!black, title=ERREURS FR\'EQUENTES,
  fonttitle=\bfseries, breakable, boxrule=0.6pt}
\newtcolorbox{fast}{colback=blue!4, colframe=blue!50!black, title=RACCOURCI PLUS RAPIDE,
  fonttitle=\bfseries, breakable, boxrule=0.6pt}
\newtcolorbox{explik}{colback=yellow!8, colframe=orange!70!black, title=COMMENT RAISONNER FACE \`A CETTE QUESTION,
  fonttitle=\bfseries, breakable, boxrule=0.6pt}

\newcommand{\exhead}[1]{\subsection{#1}}

\begin{document}

\begin{center}
  {\huge\bfseries\color{bcg} Support client}\\[4pt]
  {\large Comment raisonner sur Q1 a Q4 -- meme methode que le vrai test BCG X}
\end{center}
\vspace{4pt}\hrule\vspace{10pt}

\section{1. Objectif de ce document}
Ce PDF explique le \textbf{raisonnement}, pas seulement le code, pour
retrouver les 4 questions du vrai test BCG X CodeSignal (Data analysis,
Data collecting, Data processing, Classification) transposees sur le sujet support client. La mecanique
et les pieges sont identiques au vrai test -- seul le vocabulaire change.
Lisez ce document \textbf{apres} avoir essaye le notebook
\texttt{notebook/customer\_support\_exercices.ipynb} sans
aide, chrono en main.

\section{2. Question 1 -- Data analysis}

\exhead{Ce qu'on vous demande vraiment}
Trois \og insights \fg{} independants a calculer, puis a ranger dans un
CSV a deux colonnes (\texttt{insight\_type}, \texttt{value}). C'est le
format canonique de la Q1 : peu importe le nombre d'insights, la structure
de sortie ne change jamais.

\begin{lstlisting}
agent = pd.read_csv(f"{DATA}/agents.csv")
ticket = pd.concat([pd.read_csv(f"{DATA}/tickets_{i}.csv") for i in range(1, 5)], ignore_index=True)

average_rating = agent["rating"].mean()
pct_second_language = (agent["second_language"] != "no").mean() * 100
success_rate = (ticket["status"] == "Success").mean() * 100
\end{lstlisting}

\begin{explik}
D\'ecoupez toujours en trois reflexes independants :
\begin{itemize}[nosep]
  \item \textbf{moyenne simple} -> \texttt{.mean()} sur la colonne
        demandee, rien de plus ;
  \item \textbf{pourcentage d'une condition} -> le reflexe
        \texttt{condition.mean() * 100} (la moyenne d'un bool\'een est
        exactement sa proportion de \texttt{True}) ;
  \item \textbf{combiner plusieurs fichiers avant de calculer} ->
        \texttt{pd.concat([...], ignore\_index=True)}. Piege classique :
        calculer un taux fichier par fichier puis faire la moyenne des 4
        taux -- \textbf{faux} si les fichiers n'ont pas le meme nombre de
        lignes (moyenne non ponderee). Toujours concatener d'abord.
\end{itemize}
\end{explik}

\begin{pat}
\emph{\og Data analysis, plusieurs insights, CSV insight\_type/value \fg}
$\to$ calculer chaque insight separement sur les donnees deja
nettoyees/combinees, assembler en DataFrame a la fin,
\texttt{to\_csv(index=False)}.
\end{pat}

\begin{trap}
\begin{itemize}[nosep]
  \item Oublier de combiner les 4 fichiers \texttt{tickets\_i.csv} avant
        de calculer le taux de succes.
  \item Confondre pourcentage (\texttt{* 100}) et proportion brute --
        relisez toujours l'\'enonc\'e pour savoir laquelle est attendue.
  \item Tol\'erance de l'\'evaluateur : 2 d\'ecimales -- inutile
        d'arrondir vous-meme dans le fichier, laissez la valeur compl\`ete.
\end{itemize}
\end{trap}

\section{3. Question 2 -- Data collecting}

\exhead{Ce qu'on vous demande vraiment}
Ce n'est PAS un exercice de nettoyage : c'est un exercice de
\textbf{jointure + agr\'egation + d\'erivation de colonnes}. Vous recevez
agents.csv, headsets.csv et 4 fichiers
d'tickets, et vous devez reconstruire une seule table
avec un sch\'ema de sortie pr\'ecis.

\begin{explik}
La m\'ethode g\'en\'erale, valable pour toute question de ce style :
\begin{enumerate}[nosep]
  \item \textbf{Repartez du sch\'ema de sortie demand\'e}, colonne par
        colonne, et identifiez sa source : vient-elle directement d'une
        table, d'une jointure, d'un calcul simple, ou d'une agr\'egation
        sur une table d'\'ev\'enements ?
  \item \textbf{Sommer plusieurs colonnes bool\'eennes en une ligne} ->
        \texttt{df[bool\_cols].sum(axis=1)} donne le nombre de
        \texttt{True} par ligne ; agr\'egez ensuite par agent avec
        \texttt{groupby().sum()}.
  \item \textbf{"Nombre de jours depuis une date"} -> toujours
        \texttt{pd.to\_datetime} puis soustraction d'une date de r\'ef\'erence
        fournie par l'\'enonc\'e (ici 2023-04-15, jamais
        \texttt{pd.Timestamp.now()}), puis \texttt{.dt.days}.
  \item \textbf{Une formule explicite donn\'ee dans l'\'enonc\'e}
        (\texttt{experience = 2023 - started\_year}) doit \^etre suivie
        \`a la lettre, m\^eme si une autre m\'ethode semble \'equivalente.
  \item \textbf{fillna(0) apr\`es un merge sur une agr\'egation} -- un
        agent sans aucun \'ev\'enement n'a pas de ligne dans la
        table d'\'ev\'enements, donc le merge produit un \texttt{NaN} ; 0
        \'ev\'enement = 0 upvote, logiquement.
\end{enumerate}
\end{explik}

\begin{lstlisting}
upvote_cols = ['clarity_upvote_given', 'politeness_upvote_given', 'resolution_speed_upvote_given', 'follow_up_upvote_given']
ticket["n_upvotes"] = ticket[upvote_cols].sum(axis=1)
upvotes_per_agent = ticket.groupby("agent\_id")["n_upvotes"].sum()

headsets["last_inspection_date"] = pd.to_datetime(headsets["last_inspection_date"])
headsets["days_since_inspection"] = (TODAY - headsets["last_inspection_date"]).dt.days
\end{lstlisting}

\begin{pat}
\emph{\og table de sortie avec colonnes venant de plusieurs fichiers \fg}
$\to$ partir du sch\'ema cible, mapper chaque colonne \`a sa source
(directe / jointure / d\'eriv\'ee / agr\'eg\'ee), assembler avec des
\texttt{merge} successifs.
\end{pat}

\begin{trap}
\begin{itemize}[nosep]
  \item Les tests sont \textbf{order-agnostic} (ordre des lignes/colonnes
        indiff\'erent) -- ne perdez pas de temps \`a trier, l'\'enonc\'e le
        dit explicitement.
  \item Oublier \texttt{fillna(0)} sur \texttt{number\_of\_upvotes} apr\`es
        le merge -- les agents sans \'ev\'enement se retrouvent avec
        un \texttt{NaN} au lieu de 0.
  \item Recalculer \texttt{experience} autrement que par la formule
        exacte donn\'ee (\texttt{2023 - started\_year}).
\end{itemize}
\end{trap}

\section{4. Question 3 -- Data processing}

\exhead{La r\`egle qui domine tout : pas de fuite de donn\'ees}
Toute statistique apprise sur les donn\'ees (moyenne de remplissage,
encodeur, scaler) doit \^etre calcul\'ee \textbf{uniquement sur train},
puis appliqu\'ee telle quelle \`a test. Ne jamais recalculer quoi que ce
soit sur test.

\begin{explik}
Quatre \'etapes, quatre pi\`eges pr\'ecis -- traitez-les comme des blocs
ind\'ependants :
\begin{itemize}[nosep]
  \item \textbf{a. fillna(age, moyenne)} -- la moyenne vient du
        \textbf{train uniquement}, appliqu\'ee ensuite comme constante \`a
        test.
  \item \textbf{b. encodage ordinal} -- \texttt{sklearn.OrdinalEncoder}
        garantit nativement des codes cons\'ecutifs \`a partir de 0.
        \texttt{fit} sur train seul, \texttt{transform} (jamais
        \texttt{fit\_transform}) sur test -- sinon les deux fichiers
        peuvent avoir des mappings incoh\'erents entre eux.
  \item \textbf{c. Standard Scaling} -- m\^eme r\`egle : \texttt{fit} sur
        train, \texttt{transform} sur test.
  \item \textbf{d. binaire "A class"/"B class"} -- un simple
        \texttt{.map(\{"A class": 0, "B class": 1\})}, aucune fuite
        possible ici.
\end{itemize}
\end{explik}

\begin{trap}
\textbf{Le pi\`ege des \og exactement 5 d\'ecimales \fg.} Arrondir la
valeur num\'erique (\texttt{.round(5)}) ne suffit pas : \texttt{to\_csv}
n'ajoute pas de z\'eros de padding (\texttt{1.2} reste \texttt{"1.2"}, pas
\texttt{"1.20000"}). Formatez explicitement la colonne en cha\^ine avec
\texttt{s.map(lambda x: f"\{x:.5f\}")} avant \texttt{to\_csv} -- et
uniquement sur la colonne demand\'ee : un \texttt{float\_format="\%.5f"}
global sur \texttt{to\_csv} reformaterait \textbf{toutes} les colonnes
flottantes du DataFrame, y compris \texttt{rating}, ce qui n'est pas
demand\'e.
\end{trap}

\begin{pat}
\emph{\og Data processing avec contrainte no leakage \fg} $\to$
syst\'ematiquement \texttt{fit} sur train uniquement, \texttt{transform}
sur test ; toute statistique de remplissage vient du train.
\end{pat}

\section{5. Question 4 -- Classification}

\exhead{L'objectif est asym\'etrique, pas un bon F1}
L'\'enonc\'e dit : \emph{maximiser le recall, en gardant la precision \`a
un niveau relativement \'elev\'e}, classe positive = \texttt{B class}
(\texttt{1}). Ce n'est pas \og un bon score global \fg -- c'est une
contrainte pr\'ecise \`a traduire en code.

\begin{explik}
Deux leviers \`a utiliser ensemble :
\begin{enumerate}[nosep]
  \item \textbf{\texttt{class\_weight="balanced"}} au moment du
        \texttt{fit} -- le mod\`ele p\'enalise plus fortement les faux
        n\'egatifs sur la classe positive, donc apprend mieux \`a la
        d\'etecter.
  \item \textbf{Ajuster le seuil de d\'ecision} plut\^ot que le seuil par
        d\'efaut de \texttt{predict()} (0.5 implicite). Balayez les seuils
        sur \texttt{predict\_proba}, gardez celui qui maximise le recall
        \textbf{sous une contrainte minimale de precision} -- \c{c}a
        traduit directement la consigne de l'\'enonc\'e en code.
\end{enumerate}
\'Evaluez toujours sur \texttt{val.csv} (labels connus), jamais sur
\texttt{test.csv} (labels absents dans le vrai test -- litt\'eralement
impossible d'y calculer precision/recall en local).
\end{explik}

\begin{lstlisting}
model = LogisticRegression(max_iter=1000, class_weight="balanced")
model.fit(X_tr, y_tr)

proba_va = model.predict_proba(X_va)[:, 1]
for thr in np.arange(0.1, 0.9, 0.02):
    pred = (proba_va >= thr).astype(int)
    # garder le seuil qui maximise recall sous precision >= 0.6
\end{lstlisting}

\begin{trap}
\begin{itemize}[nosep]
  \item Le bouton Run ne score que les 10 premi\`eres lignes du fichier
        test -- ce n'est pas repr\'esentatif, ne pas paniquer si le score
        Run est bas ou instable ; seul Submit compte pour le score final.
  \item Limite d'ex\'ecution : 8 secondes -- un mod\`ele trop lourd
        (grid search complet, for\^et de 1000 arbres) peut timeout.
  \item Utiliser \texttt{predict()} par d\'efaut (seuil 0.5 implicite) au
        lieu de chercher le seuil qui correspond vraiment \`a la consigne.
\end{itemize}
\end{trap}

\begin{pat}
\emph{\og maximiser recall en gardant precision \'elev\'ee, classe
positive minoritaire \fg} $\to$ \texttt{class\_weight="balanced"} +
recherche de seuil sur \texttt{predict\_proba} avec contrainte de
precision minimale, valid\'e sur le validation set.
\end{pat}

\section{6. Check-list finale}
\begin{itemize}[nosep]
  \item[$\square$] Q1 : concat\'ener les 4 fichiers AVANT de calculer un taux global
  \item[$\square$] Q1 : \texttt{condition.mean() * 100} pour un pourcentage
  \item[$\square$] Q2 : sch\'ema de sortie decompos\'e colonne par colonne avant de coder
  \item[$\square$] Q2 : \texttt{fillna(0)} apr\`es un merge sur une agr\'egation
  \item[$\square$] Q3 : fit sur train uniquement, transform sur test -- jamais l'inverse
  \item[$\square$] Q3 : format texte \`a 5 d\'ecimales exactes, colonne cibl\'ee seulement
  \item[$\square$] Q4 : \texttt{class\_weight="balanced"} + recherche de seuil
  \item[$\square$] Q4 : \'evaluer sur val, jamais sur test (labels absents)
\end{itemize}

\end{document}
